% Actual multiplicity L-symbol construction from exact action trees.
% All six precursor modules, including the full action comparison, have elaborated.
% The fixed-final L extraction and its regression checks remain pending.
\section{Comparing exact decomposition coordinates}
\label{sec:glm_boundary_comparison}

\begin{theorem}[The reconstructed block-algebra map is unique]
    \label{thm:glm_boundary_representation_unique}
    \lean{MPSTensor.IsBiorthogonalDecomposition.sum_sandwich_eq_of_wordTupleSpanTop,
        MPSTensor.IsBiorthogonalDecomposition.support_eq_of_wordTupleSpanTop}
    \leanok
    \uses{def:glm_boundary_decomposition,
        def:blockwise_product_word_span}
    Let $B$ have two exact biorthogonal decompositions over a target family
    $(A_c)_c$, with repetitions indexed by maps $f:I\to C$ and $g:J\to C$.
    Suppose the simultaneous target word tuples span the product matrix
    algebra at one positive length. For every tuple of matrices $(M_c)_c$,
    \begin{align}
        \sum_{i\in I}W_iM_{f(i)}V_i
         & =\sum_{j\in J}\widetilde W_jM_{g(j)}\widetilde V_j,
        \label{eq:glm_boundary_representation_unique}          \\
        \sum_{i\in I}W_iV_i
         & =\sum_{j\in J}\widetilde W_j\widetilde V_j.
        \label{eq:glm_boundary_support_unique}
    \end{align}
\end{theorem}
\begin{proof}
    \leanok
    \uses{thm:glm_boundary_decomposition}
    At a simultaneous word tuple, both sides of
    (\ref{eq:glm_boundary_representation_unique}) equal the corresponding
    nonempty word of $B$. Linearity and spanning give equality on every
    tuple. Taking $M_c=I$ gives (\ref{eq:glm_boundary_support_unique}).
\end{proof}

\begin{theorem}[Coordinate changes on a common support]
    \label{thm:glm_boundary_support_comparison}
    \lean{Matrix.comparison_of_common_support}
    \leanok
    Suppose $HS=I$, $\widetilde H\widetilde S=I$, and
    $SH=\widetilde S\widetilde H$. Then
    $C=\widetilde HS$ and $D=H\widetilde S$ obey
    \begin{align}
        CD            & =I, & DC & =I,            \\
        \widetilde SC & =S, & SD & =\widetilde S.
        \label{eq:glm_boundary_support_comparison}
    \end{align}
    The common support need not be the ambient identity or self-adjoint.
\end{theorem}
\begin{proof}
    \leanok
    Substitute $SH=\widetilde S\widetilde H$ in
    $\widetilde S\widetilde HS$ and $SH\widetilde S$, then cancel the
    coordinate identities. Left multiplication by the corresponding
    analysis matrix gives the two inverse identities.
\end{proof}

\begin{theorem}[Composition and action of exact decompositions]
    \label{thm:glm_boundary_decomposition_operations}
    \lean{MPSTensor.IsBiorthogonalDecomposition.comp,
        MPSTensor.IsBiorthogonalDecomposition.sandwich,
        MPSTensor.IsBiorthogonalDecomposition.actTensor_idKron,
        MPSTensor.IsBiorthogonalDecomposition.actTensor_kronId}
    \leanok
    \uses{def:glm_boundary_decomposition}
    Exact biorthogonal decompositions compose along finite two-stage paths:
    if the intermediate and final maps are $(V_i,W_i)$ and
    $(V'_{ij},W'_{ij})$, the composite maps are
    $V'_{ij}V_i$ and $W_iW'_{ij}$. Replacing $B^a$ by $PB^aQ$ with
    $QP=I$ replaces its maps by $V_iQ$ and $PW_i$.
    Acting with a fixed MPO replaces state maps by $I\otimes V_i$ and
    $I\otimes W_i$; acting on a fixed state replaces operator maps by
    $V_i\otimes I$ and $W_i\otimes I$.
\end{theorem}
\begin{proof}
    \leanok
    Cancel the intermediate analysis-synthesis pairs in each product.
    Substitute the exact reconstruction formulas and distribute the
    finite sums. Multiplicativity of Kronecker products gives the two
    action identities.
\end{proof}

\begin{definition}[Full decomposition coordinates]
    \label{def:glm_boundary_full_coordinates}
    \lean{MPSTensor.decompositionSynthesis,MPSTensor.decompositionAnalysis}
    \leanok
    Collect the columns of all $W_i$ into $S$ and the rows of all $V_i$
    into $H$, with coordinate space $\bigoplus_i\mathbb C^{D_i}$.
\end{definition}

\begin{theorem}[Full coordinate comparison]
    \label{thm:glm_boundary_full_comparison}
    \lean{MPSTensor.decompositionSynthesis_mul_analysis,
        MPSTensor.IsBiorthogonalDecomposition.analysis_mul_synthesis,
        MPSTensor.IsBiorthogonalDecomposition.fullComparison}
    \leanok
    \uses{def:glm_boundary_full_coordinates,
        def:glm_boundary_decomposition,
        def:blockwise_product_word_span}
    The collected maps obey $HS=I$ and $SH=\sum_iW_iV_i$.
    For two exact decompositions over target blocks spanning simultaneously at one positive word length,
    $\widetilde HS$ and $H\widetilde S$ are inverse full coordinate changes,
    and each carries the other synthesis matrix to its own.
\end{theorem}
\begin{proof}
    \leanok
    \uses{thm:glm_boundary_representation_unique,
        thm:glm_boundary_support_comparison}
    Sum the separate biorthogonality identities and apply
    Theorems~\ref{thm:glm_boundary_representation_unique}
    and~\ref{thm:glm_boundary_support_comparison}.
\end{proof}

\begin{theorem}[Target intertwiners and reindexing]
    \label{thm:glm_boundary_target_intertwiners}
    \lean{MPSTensor.IsBiorthogonalDecomposition.analysis_mul_letter,
        MPSTensor.IsBiorthogonalDecomposition.letter_mul_synthesis,
        MPSTensor.IsBiorthogonalDecomposition.cross_intertwines,
        MPSTensor.IsBiorthogonalDecomposition.reindex}
    \leanok
    \uses{def:glm_boundary_decomposition}
    Exact biorthogonal decomposition implies
    $V_iB^a=A_i^aV_i$ and $B^aW_i=W_iA_i^a$.
    Consequently $V_i\widetilde W_j$ intertwines $A_i$ with
    $\widetilde A_j$. A bijective reindexing of the finite target copies
    preserves every decomposition identity.
\end{theorem}
\begin{proof}
    \leanok
    Multiply the reconstruction formula by the selected analysis or
    synthesis map and cancel all other copies. Compose the two resulting
    intertwiners. Reindex the finite reconstruction sum for the last assertion.
\end{proof}

\begin{definition}[The two exact action trees]
    \label{def:glm_boundary_action_trees}
    \lean{MPOTensor.FusionActionPath,MPOTensor.SequentialActionPath,
        MPOTensor.fusionThenActionAnalysis,MPOTensor.fusionThenActionSynthesis,
        MPOTensor.sequentialActionAnalysis,MPOTensor.sequentialActionSynthesis}
    \leanok
    \uses{def:glm_boundary_decomposition}
    Use the source's fusion analysis and synthesis $W,\widehat W$, and
    action analysis and synthesis $V,\widehat V$. The first paths are
    $(c,\mu;y,k)$ with $\mu<N_{ab}^c$ and $k<M_{cx}^y$; the second are
    $(z,j;y,i)$ with $j<M_{bx}^z$ and $i<M_{az}^y$.
    Let $\alpha$ identify the incoming right bracketing with the left.
    The tree maps are
    \begin{align}
        H_{\rm fus}^{c\mu;yk} & =V_{cx}^{y,k}(W_{ab}^{c,\mu}\otimes I),                      \\
        S_{\rm fus}^{c\mu;yk} & =(\widehat W_{ab}^{c,\mu}\otimes I)\widehat V_{cx}^{y,k},
        \label{eq:glm_boundary_fusion_action_maps}                                           \\
        H_{\rm seq}^{zj;yi}   & =V_{az}^{y,i}(I\otimes V_{bx}^{z,j})\alpha^{-1},             \\
        S_{\rm seq}^{zj;yi}   & =\alpha(I\otimes\widehat V_{bx}^{z,j})\widehat V_{az}^{y,i}.
        \label{eq:glm_boundary_sequential_action_maps}
    \end{align}
\end{definition}

\begin{theorem}[The two trees reconstruct the same action]
    \label{thm:glm_boundary_exact_action_trees}
    \lean{MPOTensor.fusionThenAction_isBiorthogonal,
        MPOTensor.sequentialAction_isBiorthogonal,
        MPOTensor.fusionThenAction_support_eq_sequentialAction}
    \leanok
    \uses{def:glm_boundary_action_trees,
        def:blockwise_product_word_span}
    If the pairwise fusion and action maps are exact and biorthogonal,
    both trees give exact biorthogonal decompositions of $(T_aT_b)\cdot A_x$.
    Their ambient support idempotents agree whenever the final state blocks
    have a simultaneous positive word span.
\end{theorem}
\begin{proof}
    \leanok
    \uses{thm:glm_boundary_decomposition_operations,
        thm:glm_boundary_representation_unique}
    Lift each first-stage decomposition through the other action factor,
    then compose with the second-stage decomposition. The bond associator
    puts both trees on the same incoming space. Apply
    Theorem~\ref{thm:glm_boundary_representation_unique} to their supports.
\end{proof}

\begin{definition}[The full action comparison]
    \label{def:glm_boundary_action_comparison}
    \lean{MPOTensor.FusionActionCoordinate,MPOTensor.SequentialActionCoordinate,
        MPOTensor.fusionToSequentialComparison,MPOTensor.sequentialToFusionComparison}
    \leanok
    \uses{def:glm_boundary_action_trees,
        def:glm_boundary_full_coordinates}
    Retaining the final bond coordinate, set
    $C=H_{\rm seq}S_{\rm fus}$ and $D=H_{\rm fus}S_{\rm seq}$.
    The source's printed L orientation is that of $C$.
\end{definition}

\begin{theorem}[Inverse full action comparisons]
    \label{thm:glm_boundary_action_comparison}
    \lean{MPOTensor.fullActionComparison_spec,
        MPOTensor.fusionToSequentialComparison_mul_analysis}
    \leanok
    \uses{def:glm_boundary_action_comparison,
        def:blockwise_product_word_span}
    Under the exact-tree and spanning assumptions,
    \begin{align}
        CD           & =I,           & DC           & =I,           \\
        S_{\rm seq}C & =S_{\rm fus}, & S_{\rm fus}D & =S_{\rm seq}, \\
        CH_{\rm fus} & =H_{\rm seq}.
        \label{eq:glm_boundary_action_comparison}
    \end{align}
\end{theorem}
\begin{proof}
    \leanok
    \uses{thm:glm_boundary_exact_action_trees,
        thm:glm_boundary_full_comparison}
    The common-support comparison gives both inverses and both synthesis
    equations. For the analysis equation, substitute
    $S_{\rm fus}H_{\rm fus}=S_{\rm seq}H_{\rm seq}$ into
    $H_{\rm seq}S_{\rm fus}H_{\rm fus}$ and cancel
    $H_{\rm seq}S_{\rm seq}=I$.
\end{proof}

\begin{definition}[Actual multiplicity L matrices]
    \label{def:glm_boundary_action_l_matrix}
    \lean{MPOTensor.FusionActionMultiplicity,MPOTensor.SequentialActionMultiplicity,
        MPOTensor.fusionActionPathEquiv,MPOTensor.sequentialActionPathEquiv,
        MPOTensor.fusionActionCoordinateEquiv,MPOTensor.sequentialActionCoordinateEquiv,
        MPOTensor.actionLMatrix,MPOTensor.inverseActionLMatrix}
    \uses{def:glm_boundary_action_comparison}
    For fixed $a,b,x,y$, the row and column spaces of $L$ have indices
    $(z,i,j)$ and $(c,k,\mu)$, respectively. Group the full coordinates
    by the final state label and define
    \begin{align}
        L_{abx}^{y}[z,i,j;c,k,\mu]
         & =D_y^{-1}\tr\bigl(H_{\rm seq}^{zj;yi}S_{\rm fus}^{c\mu;yk}\bigr).
        \label{eq:glm_boundary_action_l_entries}
    \end{align}
    Define $K_{abx}^y$ by the normalized trace of the reverse contraction.
    Individual entries may vanish.
\end{definition}

\begin{theorem}[Separation of the final bond factor]
    \label{thm:glm_boundary_action_l_factor}
    \lean{MPOTensor.actionLMatrix_cross,MPOTensor.inverseActionLMatrix_cross,
        MPOTensor.fusionToSequentialComparison_eq_blockDiagonal,
        MPOTensor.sequentialToFusionComparison_eq_blockDiagonal}
    \uses{def:glm_boundary_action_l_matrix,
        def:blockwise_product_word_span}
    Assume the pairwise fusion and action maps are exact and biorthogonal.
    Suppose the state blocks are normal and positive-dimensional, and
    their word tuples span the product algebra at one length. In grouped
    final-block coordinates,
    \begin{align}
        C & =\bigoplus_y(L_{abx}^y\otimes I_{D_y}), &
        D & =\bigoplus_y(K_{abx}^y\otimes I_{D_y}).
        \label{eq:glm_boundary_action_l_factor}
    \end{align}
\end{theorem}
\begin{proof}
    \uses{thm:glm_boundary_exact_action_trees,
        thm:glm_boundary_target_intertwiners,
        thm:glm_boundary_multiplicity_gauge}
    Each cross contraction intertwines its final blocks. Distinct final
    blocks are separated by the simultaneous word span. Within one normal
    block, the commutant is scalar; its normalized trace gives
    (\ref{eq:glm_boundary_action_l_entries}).
\end{proof}

\begin{theorem}[Invertible L matrices and their defining action equation]
    \label{thm:glm_boundary_action_l_exists}
    \lean{MPOTensor.actionLMatrix_inverse,
        MPOTensor.actionLMatrix_inverse_of_isInjective,
        MPOTensor.actionLMatrix_analysis}
    \uses{def:glm_boundary_action_l_matrix}
    Assume exact biorthogonal fusion and action maps are given, and the
    state blocks are injective, positive-dimensional, and pairwise
    inequivalent under a virtual gauge and a nonzero scalar. The matrices
    constructed from these maps satisfy
    \begin{align}
        L_{abx}^yK_{abx}^y & =I,                                                           & K_{abx}^yL_{abx}^y & =I, \\
        H_{\rm seq}        & =\left(\bigoplus_yL_{abx}^y\otimes I_{D_y}\right)H_{\rm fus}.
        \label{eq:glm_boundary_action_l_defining}
    \end{align}
    Zero multiplicity spaces are allowed. In the source index order,
    (\ref{eq:glm_boundary_action_l_defining}) is the analysis-tree relation
    of \cite[Section~3]{GarreRubio2023MPOSym}. Read the following trees from
    right to left: red legs have dimensions $\chi_a,\chi_b,\chi_c$ and
    black legs have dimensions $D_x,D_y,D_z$. The node indices are
    multiplicity labels, with row order $(z,i,j)$ and column order $(c,k,\mu)$.
    Both sides map $\C^{\chi_a}\otimes\C^{\chi_b}\otimes\C^{D_x}$ to
    $\C^{D_y}$, using the canonical identification of bracketings.
    \begin{align}
        % Source: GLM23 v3, eq:F_symbol2 (REsubmission.tex:517--534).
        % Formula: H_seq[z,i,j] = sum_{c,k,mu} L[(z,i,j),(c,k,mu)] H_fus[c,k,mu].
        % Ink/index map: red bonds a,b,c are operator virtual spaces chi_a,chi_b,chi_c;
        % black bonds x,y,z are state virtual spaces D_x,D_y,D_z. All ports are virtual.
        % Main-text V is action analysis; W is fusion analysis. No hats or adjoints.
        % Sequential contraction: V_az^{y,i} (I_chi_a tensor V_bx^{z,j}).
        % Fusion contraction: V_cx^{y,k} (W_ab^{c,mu} tensor I_Dx).
        % Only z is contracted in the left tree; only c is contracted in each right tree.
        % Open legs on both sides: final y on the west; inputs a,b,x top-to-bottom east.
        % Boundary signature on each side: (virtual-west=1 | virtual-east=3 |
        % physical-up=0 | physical-down=0). Parenthesization is identified canonically.
        % The sum over c,k,mu is scalar linear combination, not extra bond contraction.
        \begingroup
        \tenkzkernel
        \tndeclare{species}{glm-operator}{hue=source:red!70!black}
        \tndeclare{species}{glm-state}{hue=source:black}
        \begin{tenkzeq}[check=signature]
            \begin{tenkz}[rows=3,cols=2,bonds=none]
                \tn[at=(1,1),name=outer,skin=box,wires=2,
                    ports={180:virtual:$y$,0@1:virtual:$a$,0@2:virtual}]{V^{i}}
                \tn[at=(2,2),name=inner,skin=box,wires=2,
                    ports={180:virtual,0@1:virtual:$b$,0@2:virtual:$x$}]{V^{j}}
                \tn[at=(1,2),void=sealed,skin=none]{}
                \tnwire[species=glm-state]{outer.180}{open w}
                \tnwire[species=glm-operator]{outer.0@1}{open e}
                \tnwire[species=glm-state,route=arc,name=middle-state]{outer.0@2}{inner.180}
                \tnmark[form=label,label pos=90]{on middle-state 0.5}{$z$}
                \tnwire[species=glm-operator]{inner.0@1}{open e}
                \tnwire[species=glm-state]{inner.0@2}{open e}
            \end{tenkz}
            = \sum_{c,k,\mu} (L_{abx}^{y})_{c,k\mu}^{z,ij}
            \begin{tenkz}[rows=3,cols=2,bonds=none]
                \tn[at=(2,1),name=action,skin=box,wires=2,
                    ports={180:virtual:$y$,0@1:virtual,0@2:virtual:$x$}]{V^{k}}
                \tn[at=(1,2),name=fusion,skin=box,wires=2,
                    ports={180:virtual,0@1:virtual:$a$,0@2:virtual:$b$}]{W^{\mu}}
                \tn[at=(1,1),void=sealed,skin=none]{}
                \tnwire[species=glm-state]{action.180}{open w}
                \tnwire[species=glm-operator,route=arc,name=middle-operator]{action.0@1}{fusion.180}
                \tnmark[form=label,label pos=270]{on middle-operator 0.5}{$c$}
                \tnwire[species=glm-state]{action.0@2}{open e}
                \tnwire[species=glm-operator]{fusion.0@1}{open e}
                \tnwire[species=glm-operator]{fusion.0@2}{open e}
            \end{tenkz}
        \end{tenkzeq}
        \endgroup
        \label{eq:glm_boundary_action_l_diagram}
    \end{align}
\end{theorem}
\begin{proof}
    \uses{thm:glm_boundary_action_l_factor,
        thm:glm_boundary_action_comparison,
        thm:glm_boundary_spanning}
    Derive a common positive word span. Substitute
    (\ref{eq:glm_boundary_action_l_factor}) in the full inverse and analysis
    equations. Take each final block and cancel its nonzero identity
    factor to obtain (\ref{eq:glm_boundary_action_l_defining}).
\end{proof}

\paragraph{Source index convention.}
The coupled pentagon printed in the source reverses both multiplicity pairs
in its fusion factor. The typed correction and the analysis-gauge convention
are recorded in \cite{gap:glm23_multiplicity_l_indices}. The present construction
proves the invertible action comparison; the corrected coupled pentagon and
its general multiplicity-gauge law remain separate proof obligations.
